\section{Введение}\label{sec:intro}
Стохастическая аппроксимация с одновременным возмущением (SPSA) оценивает градиент по двум измерениям
целевой функции независимо от размерности задачи \cite{spall1992}.
Поэтому метод применяют в задачах, где градиент недоступен, а измерения зашумлены.
Скорость сходимости SPSA определяется шагами $a_k$ и $c_k$, и при сильном шуме траектория
колеблется вокруг оптимума.
В работе исследуется добавление момента к оценке градиента.
Мы показываем, что момент $\beta = 0.6$ уменьшает итоговую ошибку на зашумлённой квадратичной функции
размерности 10 при том же числе измерений.
